\documentclass[spanish,letterpaper,oneside]{article}
\def\documenttitle {Delitos contra el patrimonio}
\def\documentsubtitle {An\'alisis comparativo de figuras y subsunci\'on de hechos}
\def\documentsubject {Actividad S7 - An\'alisis de hechos}
\def\documentauthor {Mart\'in Jonathan de la Cruz}
\def\coursename {Derecho penal especial mexicano}
\def\coursecode {020}
\def\hfwidthwrap {true}
\def\hfwidthtitle {0.52}
\def\hfwidthcourse {0.42}
\def\universityname {Universidad Abierta y a Distancia de M\'exico}
\def\universityfaculty {Licenciatura en Derecho}
\def\universitydepartment {Derecho penal especial mexicano}
\def\universitydepartmentimage {departamentos/UnADM}
\def\universitydepartmentimagecfg {height=1.57cm}
\def\universitylocation {Ciudad de M\'exico}
\def\authortable {
  \begin{tabular}{ll}
    \textbf{Alumno:} & Mart\'in Jonathan de la Cruz \\
    Matr\'icula: & ES2611202040 \\
    Figura docente: & Ver\'onica Ortiz L\'opez \\
    Grupo: & 020 \\
    Semestre/Bloque: & 2 / 2 \\
    Semana: & 7 \\
    \multicolumn{2}{l}{Fecha: 15 de noviembre de 2026} \\
    \multicolumn{2}{l}{\universitylocation}
  \end{tabular}
}

% Wrapper local para apuntar a la plantilla maestra de AulaTeX
\input{C:/Users/Sysx/Documents/AulaTeX-Academico/aulaTEX/base/Plantilla-Informe/template.tex}

\def\indexstyle{}
\def\nameltcont{Tabla de contenido}
\setcounter{tocdepth}{2}
\usepackage{helvet}
\renewcommand{\familydefault}{\sfdefault}
\usepackage[most]{tcolorbox}
\usepackage{array,tabularx,booktabs,longtable}
\usepackage{tikz}
\usetikzlibrary{babel}
\definecolor{casoverde}{HTML}{174A3A}
\definecolor{casopapel}{HTML}{F6F7F2}
\definecolor{tableHeader}{HTML}{E8EFE5}
\definecolor{rowAlt}{HTML}{F7FAF7}
\setcitestyle{authoryear,open={(},close={)},citesep={;}}
\def\pagemargintop{2.2}
\def\pagemarginbottom{2.2}
\def\predocpageromannumber{true}
\def\predocresetpagenumber{true}

\newcolumntype{L}[1]{>{\raggedright\arraybackslash}p{#1}}

\newcommand{\insertcoverwatermark}{%
  \AddToShipoutPictureBG*{%
    \begin{tikzpicture}[remember picture,overlay]
      \node[opacity=0.13,inner sep=0pt,yshift=-0.35cm]
        at (current page.center) {\includegraphics[width=0.72\paperwidth]{img/departamentos/UnADM.pdf}};
    \end{tikzpicture}%
  }%
}

\begin{document}
\insertcoverwatermark
\templatePortrait
\templatePagecfg
\onehalfspacing

\begin{abstractd}
El estudio dogm\'atico de los delitos contra el patrimonio exige un deslinde anal\'itico riguroso de sus elementos t\'ipicos objetivos, normativos y subjetivos, a fin de garantizar el principio de taxatividad y exacta aplicaci\'on de la ley penal (art\'iculo 14 de la Constituci\'on Pol\'itica de los Estados Unidos Mexicanos). Las figuras de robo, abuso de confianza, fraude, despojo y da\~no en propiedad ajena tutelan el acervo econ\'omico de las personas, pero difieren sustancialmente en la modalidad de afectaci\'on, el t\'itulo originario de tenencia del bien y el despliegue conductual del agente activo. El presente reporte desarrolla una investigaci\'on jur\'idica exhaustiva sustentada en el C\'odigo Penal Federal (CPF) y en el C\'odigo Penal para el Distrito Federal (CPDF), integrando una tabla comparativa con la delimitaci\'on descriptiva y el rasgo diferenciador de cada figura, acompa\~nada de la fundamentaci\'on de una circunstancia agravante del robo. Posteriormente, se realiza la subsunci\'on preliminar de cinco hechos f\'acticos agrupados en tres fichas, identificando con precisi\'on los elementos normativos acreditados, la informaci\'on f\'altante para una calificaci\'on definitiva y la hip\'otesis legal aplicable, concluyendo con una s\'intesis dogm\'atica en estricto apego a la extensi\'on m\'axima de siete cuartillas.\footnote{Declaraci\'on sobre el uso \'etico de inteligencia artificial: En cumplimiento de los lineamientos institucionales de la UnADM, se declara el uso instrumental de un modelo de lenguaje artificial \'unicamente para la revisi\'on de estilo y maquetaci\'on editorial. El razonamiento jur\'idico, la clasificaci\'on dogm\'atica y las conclusiones son de autor\'ia propia y exclusiva del estudiante.}
\end{abstractd}

\templateIndex
\templateFinalcfg

\section{Investigaci\'on jur\'idica: Delitos contra el patrimonio}

El patrimonio econ\'omico constituye un bien jur\'idico medular tutelado por el orden penal, concebido como el conjunto de derechos, bienes y obligaciones susceptibles de apreciaci\'on pecuniaria que pertenecen leg\'itimamente a una persona f\'isica o moral \citep{amuchateguiPenal2012,pavonVasconcelos2018}. Sin embargo, la lesi\'on patrimonial no opera de manera un\'ivoca; el legislador mexicano ha tipificado modalidades diferenciadas seg\'un el modo de acceso al objeto material y la posici\'on de la v\'ictima.

\subsection{Tabla de investigaci\'on de figuras patrimoniales}

En el cuadro \ref{tab:investigacion_patrimonial} se sintetizan las cinco figuras jur\'idicas requeridas, contrastando su fundamentaci\'on en el C\'odigo Penal Federal \citep{cpfOficial2026} y en el C\'odigo Penal para el Distrito Federal \citep{cpdfOficial2026}.

\begin{table}[h!]
\centering
\caption{Delimitaci\'on dogm\'atica de las figuras t\'ipicas contra el patrimonio}
\label{tab:investigacion_patrimonial}
\scriptsize
\begin{tabularx}{\textwidth}{p{2.2cm} p{2.8cm} X p{3.6cm}}
\toprule
\rowcolor{tableHeader}
\textbf{Figura} & \textbf{Ley y art\'iculo} & \textbf{Conducta descrita con palabras propias} & \textbf{Dato que permite diferenciarla} \\
\midrule
\textbf{Robo} & 
CPF: Arts. 367, 368.\newline
CPDF: Art. 220. & 
Apoderarse materialmente de una cosa mueble ajena, traslad\'andola a la propia esfera de custodia, sin derecho y sin el consentimiento de quien legalmente puede disponer de ella. & 
El sujeto activo \textbf{no tiene previamente la tenencia ni custodia} de la cosa; rompe la esfera de custodia ajena sin t\'itulo jur\'idico previo. \\
\midrule
\rowcolor{rowAlt}
\textbf{Abuso de confianza} & 
CPF: Art. 382.\newline
CPDF: Art. 227. & 
Disponer indebidamente de un bien mueble ajeno, en beneficio propio o de un tercero, habiendo recibido previamente s\'olo su tenencia material precaria (custodia, dep\'osito o pr\'estamo de uso), sin traslaci\'on de dominio. & 
El sujeto activo \textbf{recibe l\'icitamente la tenencia precaria de buena fe} entregada por la v\'ictima; el delito surge cuando traiciona la confianza disponiendo del bien como due\~no. \\
\midrule
\textbf{Fraude} & 
CPF: Art. 386.\newline
CPDF: Art. 230. & 
Obtener un lucro econ\'omico indebido para s\'i o para otro, enga\~nando a la v\'ictima mediante maquinaciones falaces o aprovech\'andose del error en que esta se halla. & 
El desplazamiento patrimonial es \textbf{entregado voluntariamente por la v\'ictima}, pero su consentimiento se encuentra viciado por el enga\~no o maquinaci\'on previa del activo. \\
\midrule
\rowcolor{rowAlt}
\textbf{Despojo} & 
CPF: Art. 395.\newline
CPDF: Art. 237. & 
Ocupar materialmente un inmueble ajeno o hacer uso de \'el o de un derecho real que no le pertenezca, de propia autoridad y mediante violencia, enga\~no o furtividad, desposeyendo a su leg\'itimo poseedor. & 
Recae estrictamente sobre \textbf{bienes inmuebles o derechos reales} (no muebles), vulnerando la posesi\'on material o pac\'ifica a trav\'es de v\'ias de hecho no consentidas. \\
\midrule
\textbf{Da\~no en propiedad ajena} & 
CPF: Arts. 397, 399.\newline
CPDF: Art. 239. & 
Destruir, deteriorar, quebrar o inutilizar un bien mueble o inmueble ajeno, por cualquier medio, mermando su integridad f\'isica o su valor econ\'omico. & 
El m\'ovil delictivo es la \textbf{destrucci\'on o deterioro material} del bien, careciendo de \'animo de lucro, apropiaci\'on o transmisi\'on de posesi\'on para s\'i. \\
\bottomrule
\end{tabularx}
\end{table}

\subsection{Circunstancia agravante del robo y fundamento legal}

\begin{tcolorbox}[colback=casopapel,colframe=casoverde,arc=1mm,left=3mm,right=3mm,top=2mm,bottom=2mm]
\textbf{Circunstancia agravante investigada:} Comisi\'on del robo respecto de veh\'iculo automotriz o sus partes, o en v\'ia p\'ublica.\newline
\textbf{Fundamentaci\'on legal:} 
\begin{itemize}
  \item \textbf{C\'odigo Penal para el Distrito Federal (CPDF):} Art\'iculo 224, fracci\'on IX (cuando se cometa en v\'ia p\'ublica o respecto de veh\'iculo automotriz o partes de este) y fracci\'on III (por medio de la violencia f\'isica o moral).
  \item \textbf{C\'odigo Penal Federal (CPF):} Art\'iculo 381, fracci\'on VIII (cuando se cometa sobre veh\'iculo o en paraje solitario/lugar p\'ublico concurrido) y Art\'iculo 381 Bis (robo de veh\'iculo automotor).
\end{itemize}
\textbf{Criterio jurisprudencial:} Conforme a la jurisprudencia de la Suprema Corte de Justicia de la Naci\'on \citep{scjnRoboCalificado}, las calificativas no se presumen; exigen la comprobaci\'on plena de los datos objetivos normativos descritos en el tipo penal complementado, tales como el espacio f\'isico p\'ublico y la fractura de dispositivos de seguridad.
\end{tcolorbox}

\section{An\'alisis comparativo de los cinco hechos (Fichas 1 a 3)}

La subsunci\'on penal exige confrontar los hechos estrictamente acreditados con los elementos t\'ipicos de la norma, distinguiendo con claridad entre datos comprobados, deducciones l\'ogicas y elementos pendientes de probanza \citep{pavonVasconcelos2018}. En el cuadro \ref{tab:analisis_hechos} se analiza la totalidad de los hechos contenidos en las tres fichas.

\begin{table}[h!]
\centering
\caption{Matriz de subsunci\'on t\'ipica y an\'alisis de elementos normativos}
\label{tab:analisis_hechos}
\scriptsize
\begin{tabularx}{\textwidth}{p{2.3cm} p{3.2cm} p{3.4cm} p{3.8cm}}
\toprule
\rowcolor{tableHeader}
\textbf{Hechos conocidos} & \textbf{Elementos legales identificados} & \textbf{Datos que faltan} & \textbf{Posible figura y fundamento legal} \\
\midrule
\textbf{Hecho 1: Bicicleta}\newline
(Ficha 1)\newline
Una persona toma sin autorizaci\'on una bicicleta ajena que estaba estacionada. & 
$\bullet$ Sujeto activo toma (apoderamiento material) una cosa mueble ajena (bicicleta).\newline
$\bullet$ Ausencia de consentimiento de quien leg\'itimamente puede disponer del bien.\newline
$\bullet$ Falta de derecho en la conducta. & 
\textbf{Para valorar agravante se requerir\'ia acreditar:} si estaba asegurada con candado o cadena (fractura/violencia sobre las cosas), el lugar exacto de estacionamiento (v\'ia p\'ublica, paraje solitario o inmueble habitado) y la hora (nocturnidad). Tambi\'en falta el \emph{valor comercial} del bien para fijar la penalidad aplicable. & 
\textbf{Robo simple}\newline
CPF: Arts. 367, 368 y 370.\newline
CPDF: Art. 220.\newline
\emph{(No procede agravar sin constancia objetiva del medio o lugar de comisi\'on)}. \\
\midrule
\rowcolor{rowAlt}
\textbf{Hecho 2: Computadora}\newline
(Ficha 1)\newline
Alguien recibe voluntariamente una computadora prestada y despu\'es la vende sin consentimiento. & 
$\bullet$ Entrega voluntaria y leg\'itima de la tenencia material (pr\'estamo de uso/comodato).\newline
$\bullet$ Objeto material: bien mueble ajeno.\newline
$\bullet$ Acto posterior de disposici\'on ileg\'itima mediante venta no autorizada en perjuicio de la propietaria. & 
Falta documentar el t\'itulo jur\'idico formal de transmisi\'on de tenencia (pr\'estamo verbal o escrito), la acreditaci\'on formal de la venta a un tercero y el aval\'uo comercial de la computadora para la cuantificaci\'on exacta de la pena. & 
\textbf{Abuso de confianza}\newline
CPF: Art. 382.\newline
CPDF: Art. 227.\newline
\emph{(El activo ten\'ia la posesi\'on material leg\'itima; delinque al disponer del bien)}. \\
\midrule
\textbf{Hecho 3: Venta internet}\newline
(Ficha 2)\newline
Anuncia por internet art\'iculo inexistente, recibe pago mediante enga\~no y no lo entrega. & 
$\bullet$ Conducta falaz y maquinaci\'on previa (publicar art\'iculo ficticio).\newline
$\bullet$ Enga\~no determinante que induce al error a la v\'ictima.\newline
$\bullet$ Entrega voluntaria de dinero por la v\'ictima y obtenci\'on de un lucro indebido para el activo. & 
Falta identificar el importe dinerario transferido, la titularidad registral de la cuenta bancaria receptora y si el agente oper\'o de forma individual o mediante estructura tecnol\'ogica reiterada. & 
\textbf{Fraude gen\'erico}\newline
CPF: Art. 386.\newline
CPDF: Art. 230.\newline
\emph{(Desplazamiento voluntario derivado de voluntad viciada por maquinaci\'on previa)}. \\
\midrule
\rowcolor{rowAlt}
\textbf{Hecho 4: Terreno}\newline
(Ficha 2)\newline
Alguien ocupa sin derecho un terreno que otra persona posee y le impide utilizarlo. & 
$\bullet$ Inmueble ajeno determinado (terreno).\newline
$\bullet$ Acci\'on de ocupaci\'on material f\'isica de propia autoridad.\newline
$\bullet$ Turbaci\'on total que impide el uso y goce a su leg\'itimo poseedor. & 
\textbf{Dato crucial faltante:} El medio de comisi\'on empleado para la ocupaci\'on (si medi\'o violencia f\'isica, moral/amenazas, enga\~no o furtividad/clandestinidad) y el t\'itulo jur\'idico que acredita la posesi\'on previa de la v\'ictima. & 
\textbf{Despojo de inmueble}\newline
CPF: Art. 395, fracci\'on I.\newline
CPDF: Art. 237, fracci\'on I.\newline
\emph{(Tutela la posesi\'on pac\'ifica de inmuebles frente a v\'ias de hecho)}. \\
\midrule
\textbf{Hecho 5: Cristal de auto}\newline
(Ficha 3)\newline
Rompe intencionalmente el cristal del autom\'ovil ajeno en discusi\'on en lugar p\'ublico. & 
$\bullet$ Acci\'on destructiva intencional (dolo directo) provocada durante discusi\'on.\newline
$\bullet$ Objeto material: cristal de veh\'iculo automotor ajeno.\newline
$\bullet$ Menoscabo econ\'omico e inutilizaci\'on de parte del bien, sin \'animo de apropiaci\'on. & 
Falta el dictamen pericial de valuaci\'on del costo de reposici\'on del cristal; y verificar si el agresor sustrajo o intent\'o sustraer objetos del interior del veh\'iculo (lo que recalificar\'ia la conducta a tentativa de robo calificado). & 
\textbf{Da\~no a la propiedad ajena}\newline
CPF: Art. 399.\newline
CPDF: Art. 239.\newline
\emph{(Afectaci\'on destructiva al bien sin finalidad de lucro o apoderamiento)}. \\
\bottomrule
\end{tabularx}
\end{table}

\section{Conclusi\'on jur\'idica breve}

El elemento dogm\'atico m\'as determinante para deslindar las figuras patrimoniales radica en el t\'itulo y modo originario de acceso al bien: mientras en el robo se fractura la custodia ajena mediante apoderamiento material directo sin consentimiento, en el abuso de confianza la v\'ictima transmite voluntariamente la tenencia precaria de buena fe, surgiendo el il\'icito s\'olo cuando el sujeto activo dispone indebidamente del objeto como propio; en contraste, en el fraude la entrega es voluntaria pero su consentimiento est\'a viciado de origen por maquinaciones enga\~nosas. Identificar con precisi\'on este elemento de acceso inicial resulta indispensable para impedir recalificaciones err\'oneas y garantizar la estricta tipicidad penal consagrada en el art\'iculo 14 de la Constituci\'on Federal.

\clearpage
\begin{thebibliography}{99}

\bibitem[Amuchategui Requena(2012)]{amuchateguiPenal2012}
Amuchategui Requena, I. (2012). \emph{Derecho penal} (4.ª ed.). M\'exico: Oxford University Press.

\bibitem[Asamblea Legislativa del Distrito Federal(2026)]{cpdfOficial2026}
Asamblea Legislativa del Distrito Federal. (2026). \emph{C\'odigo Penal para el Distrito Federal} (Texto normativo consolidado y reformas publicadas en la Gaceta Oficial). Ciudad de M\'exico: Gobierno de la Ciudad de M\'exico.

\bibitem[C\'amara de Diputados(1917)]{cpeumAulaS3}
C\'amara de Diputados del H. Congreso de la Uni\'on. (1917). \emph{Constituci\'on Pol\'itica de los Estados Unidos Mexicanos} (Texto original con reformas publicadas en el Diario Oficial de la Federaci\'on hasta 2026). M\'exico: Secretar\'ia de Servicios Parlamentarios.

\bibitem[C\'amara de Diputados(2026)]{cpfOficial2026}
C\'amara de Diputados del H. Congreso de la Uni\'on. (2026). \emph{C\'odigo Penal Federal} (Texto vigente y reformas publicadas en el Diario Oficial de la Federaci\'on). M\'exico: Secretar\'ia de Servicios Parlamentarios.

\bibitem[Pav\'on Vasconcelos(2018)]{pavonVasconcelos2018}
Pav\'on Vasconcelos, F. (2018). \emph{Derecho penal mexicano: Parte especial (Delitos contra el patrimonio)}. M\'exico: Editorial Porr\'ua.

\bibitem[Suprema Corte de Justicia de la Naci\'on(2018)]{scjnRoboCalificado}
Suprema Corte de Justicia de la Naci\'on. (2018). \emph{Jurisprudencia 1a./J. 22/2018: Robo. Para tener por acreditadas las agravantes relativas al lugar y al medio de comisi\'on, se requiere prueba plena de los elementos objetivos del tipo complementado}. Gaceta del Semanario Judicial de la Federaci\'on, D\'ecima \'Epoca.

\bibitem[Universidad Abierta y a Distancia de M\'exico(2026)]{unadmPenalS7Plan}
Universidad Abierta y a Distancia de M\'exico. (2026). \emph{Planificaci\'on de actividades Semana 7: An\'alisis de hechos (Delitos contra el patrimonio)}. Aula virtual UnADM, Curso 3482.

\end{thebibliography}

\end{document}
